\section{Models}
\label{sec:model}

This section introduces the three building blocks of the proposed
models---the DeGroot model~\cite{DeGroot1974}, the cognitive bias
framework~\cite{Alvim2024}, and the Spiral-of-Silence opinion
models~\cite{aranda2024sound}---and defines
\CBSOMminus\ and \CBSOMplus, the Cognitive-Bias Silence Opinion
Models that combine all three.

% ===================================================================
\subsection{The DeGroot Model}
\label{sec:degroot}

The DeGroot model~\cite{DeGroot1974} formalizes social learning as
repeated weighted averaging over a fixed \emph{influence graph}.

\begin{definition}[Influence Graph]\label{def:influence-graph}
  An \emph{influence graph} is a weighted directed graph $G=(A,E,I)$
  with $A=\{1,\ldots,n\}$, $n>1$, the set of \emph{agents};
  $E\subseteq A\times A$ the set of directed edges; and
  $I:A\times A\to[0,1]$ the weight function satisfying
  $I(i,j)=0\iff(i,j)\notin E$ and, for each $i\in A$,
  \begin{equation}\label{eq:row-stochastic}
    \sum_{j\in N(i)\cup\{i\}}I(j,i)=1,
  \end{equation}
  where $N(i)=\{j\in A\setminus\{i\}:(j,i)\in E\}$ is the set of
  \emph{influencers} of~$i$.
\end{definition}

The vertices in $A$ represent the $n$ agents of a community. An edge
$(j,i)\in E$ means that agent $j$ influences agent $i$, and the weight
$I(j,i)$ denotes the strength of that influence: the higher the value,
the stronger the influence. Equation~\eqref{eq:row-stochastic} requires
that the weights that agent $i$ gives to its influencers and to itself
add up to~$1$.

Each agent $i$ holds an \emph{opinion} $B_i^t\in[0,1]$ at time $t$,
interpreted as the degree of agreement with an underlying proposition
(e.g., \emph{``AI poses a threat to humanity''}). If $B_i^t=0$, agent $i$
completely disagrees with the proposition; if $B_i^t=1$, it completely
agrees with it. The full opinion vector is $\mathbf{B}^t\in[0,1]^n$.
The synchronous DeGroot update reads:
\begin{equation}\label{eq:degroot}
  B_i^{t+1} = \sum_{j=1}^n I_{ij}\,B_j^t,
  \qquad\text{equivalently}\quad
  \mathbf{B}^{t+1}=I\,\mathbf{B}^t,
\end{equation}
where $I=[I_{ij}]$ is row-stochastic.

\begin{theorem}[DeGroot Consensus~\cite{DeGroot1974}]\label{thm:degroot}
  If $I$ is row-stochastic and primitive (i.e., $\exists\,k$ such
  that all entries of $I^k$ are positive), then there exists
  $\mathbf{B}^*$ with identical components such that
  $\lim_{t\to\infty}\mathbf{B}^t=\mathbf{B}^*$.
\end{theorem}

Despite its tractability, the DeGroot model makes two assumptions
challenged by real social dynamics: it assumes that agents always
express their opinions freely (no strategic silence), and it assumes
\emph{linear} assimilation of neighbors' opinions, ignoring nonlinear
psychological responses to disagreement~\cite{Alvim2024,NoelleNeumann1984}.
Sections~\ref{sec:biases} and~\ref{sec:som} address each limitation
in turn, before they are combined in
Sections~\ref{sec:cbsom-minus}--\ref{sec:cbsom-plus}.

% ===================================================================
\subsection{Cognitive Bias Functions}
\label{sec:biases}

Following \cite{Alvim2024}, we model cognitive bias as a function that
nonlinearly transforms the opinion disagreement between an agent and
each of its influencers before the update is applied.

\begin{definition}[Bias Function~\cite{Alvim2024}]\label{def:bias-fn}
  A \emph{bias function} for agent $i$ with respect to agent $j$ is
  a continuous map $\beta_{i,j}:[-1,1]\to[-1,1]$
  that transforms the raw disagreement $x=B_j^t-B_i^t$.
\end{definition}

Intuitively, $\beta_{i,j}(x)$ is the part of the disagreement $x$ that
agent $i$ actually absorbs from agent $j$. The identity function
$\beta_{i,j}(x)=x$ corresponds to the linear response of the DeGroot model.
The bias-extended DeGroot update replaces the linear term
$(B_j-B_i)$ with $\beta_{i,j}(B_j-B_i)$:
\begin{equation}\label{eq:degroot-bias}
  B_i^{t+1}
  = B_i^t + \sum_{j\in N(i)} I_{ji}\,
    \beta_{i,j}(B_j^t-B_i^t).
\end{equation}
The original formulation of~\cite{Alvim2024} includes an explicit
$[\,\cdot\,]_0^1$ clamp to enforce $B_i^{t+1}\in[0,1]$. Our formal results
assume bias functions in the Receptive-Resistant region $R$ defined below.
In that region the update stays within $[0,1]$ without clamping
(Lemma~\ref{lem:cbsom-bounds}), so we omit the clamp from here on.
The four representative bias functions examined in this paper are:
\begin{itemize}
  \item \textbf{Confirmation bias} (\textsl{conf}): agents attenuate
    distant disagreements,
    $\textsl{conf}(x)=x(1+\delta-|x|)/(1+\delta)$, $\delta>0$
    \cite{Alvim2024,Nyhan2010}.
  \item \textbf{Backfire effect} (\textsl{backf}): large disagreements
    trigger repulsion, $\textsl{backf}(x)=-x^3$
    \cite{Alvim2024,Friedrich2020}.
  \item \textbf{Authority bias} (\textsl{fan}): agents follow
    perceived authorities with maximal response,
    $\textsl{fan}(x)=x/|x|$ for $x\ne 0$, $0$ otherwise
    \cite{Alvim2024,Aronson2010}.
  \item \textbf{Insular bias} (\textsl{ins}): agents ignore all
    external influence, $\textsl{ins}(x)=0$ \cite{Alvim2024}.
\end{itemize}

\noindent\textbf{Bias regions.}
Alvim et al.~\cite{Alvim2024} classify bias functions by the
region of the square $[-1,1]^2$ in which the graph of $(x,\beta(x))$ lies:
\begin{itemize}
  \item \emph{Receptive-Resistant} ($R$): $0<|\beta(x)|<|x|$;
    confirmation bias belongs to this region. In synchronous DeGroot
    dynamics, consensus is preserved on strongly connected graphs
    under continuous $R$-biases~\cite{Alvim2024}.
  \item \emph{Malleability} ($M$): $|\beta(x)|\ge|x|$; authority
    bias is the extreme case.
  \item \emph{Backfire}: sign reversal,
    $\operatorname{sign}(\beta(x))\ne\operatorname{sign}(x)$;
    consensus typically fails.
  \item \emph{Insular} ($I$): zero response, $\beta(x)=0$.
\end{itemize}
The region $R$ is central to this paper. A bias in $R$ moves an agent
towards each influencer, but never beyond it. Together with continuity,
this property is the key analytical lever for extending consensus
guarantees to the models introduced below.

% ===================================================================
\subsection{Spiral-of-Silence Opinion Models}
\label{sec:som}

The Spiral of Silence~\cite{NoelleNeumann1974,NoelleNeumann1984}
posits that agents perceiving themselves as a minority tend to
self-censor to avoid social isolation.
Aranda et al.~\cite{aranda2024sound} formalize this
mechanism as an augmentation of the DeGroot framework.

\begin{definition}[Silence Opinion Model~\cite{aranda2024sound}]
\label{def:som}
  A \emph{Silence Opinion (SO) model} is a quadruple
  $(G,\mathbf{B}^0,\mathbf{S}^0,\mu_G)$, where $G=(A,E,I)$ is an
  influence graph, $\mathbf{B}^0\in[0,1]^n$ is the initial opinion
  vector, $\mathbf{S}^0\in\{0,1\}^n$ is the initial silence vector
  ($S_i^t=1$ meaning agent $i$ is \emph{vocal} at time $t$), and
  $\mu_G:[0,1]^n\times\{0,1\}^n\times\mathbb{N}
   \to[0,1]^n\times\{0,1\}^n$ is the joint update function.
\end{definition}

\begin{definition}[Consensus]\label{def:consensus}
  Agents in an SO model \emph{reach consensus} if $\exists\,v\in[0,1]$
  such that $\lim_{t\to\infty}B_i^t=v$ for all $i\in A$.
\end{definition}

\subsubsection{Memoryless Model \texorpdfstring{\SOMminus}{SOM-}}
\label{sec:som-minus}

In \SOMminus, silent agents contribute nothing to their neighbors'
updates~\cite{aranda2024sound}. Writing $x\sim_{\tau_i}y$ iff
$|x-y|\le\tau_i$ for tolerance radius $\tau_i\in[0,1]$:
\begin{align}
  B_i^{t+1} &= B_i^t
    + \sum_{j\in N(i)}I_{ji}\cdot S_j^t\cdot(B_j^t-B_i^t),
  \label{eq:som-minus-b}\\
  S_i^{t+1} &=
  \begin{cases}
    1 & \text{if }\;\bigl\lceil|N^t(i)|/2\bigr\rceil
        \le\bigl|\{j\in N^t(i)\mid B_i^t\sim_{\tau_i}B_j^t\}\bigr|,\\
    0 & \text{otherwise,}
  \end{cases}
  \label{eq:som-minus-s}
\end{align}
where $N^t(i)=\{j\in N(i):S_j^t=1\}$.
Intuitively, agent $i$ goes silent when a majority of its currently
vocal neighbors hold opinions too far from its own. Silent agents still
revise their own opinions through~\eqref{eq:som-minus-b}; they only stop
expressing them.

\subsubsection{Memory-Based Model \texorpdfstring{\SOMplus}{SOM+}}
\label{sec:som-plus}

In \SOMplus, silent agents retain influence through the last opinion
they publicly expressed~\cite{aranda2024sound}.
Assuming $\mathbf{S}^0=\mathbf{1}_n$, let
$t_j'=\max\{u\le t:S_j^u=1\}$ and
$\hat{B}_j^t=B_j^{t_j'}$ (\emph{last public opinion} of agent~$j$):
\begin{align}
  B_i^{t+1} &= B_i^t
    + \sum_{j\in N(i)}I_{ji}\cdot(\hat{B}_j^t-B_i^t),
  \label{eq:som-plus-b}\\
  S_i^{t+1} &=
  \begin{cases}
    1 & \text{if }\;\bigl\lceil|N(i)|/2\bigr\rceil
        \le\bigl|\{j\in N(i)\mid B_i^t\sim_{\tau_i}\hat{B}_j^t\}\bigr|,\\
    0 & \text{otherwise.}
  \end{cases}
  \label{eq:som-plus-s}
\end{align}
Unlike \SOMminus, the model \SOMplus\ is non-Markovian: its evolution
depends on the history of silence events through $\hat{B}_j^t$.

\subsubsection{Key Results for \texorpdfstring{\SOMminus}{SOM-}}
\label{sec:som-results}

The following results from~\cite{aranda2024sound}
underpin the analysis in Section~\ref{sec:analysis}.

\begin{lemma}[Bounded opinions~\cite{aranda2024sound}]
\label{lem:bounds}
  For every \SOMminus\ model and all $t$:
  $\min(\mathbf{B}^t)\le B_i^{t+1}\le\max(\mathbf{B}^t)$
  for all $i\in A$.
\end{lemma}

\begin{corollary}[Monotone extremes~\cite{aranda2024sound}]
\label{cor:mono}
  In every \SOMminus\ model, $\max(\mathbf{B}^t)$ is non-increasing
  and $\min(\mathbf{B}^t)$ is non-decreasing in~$t$.
\end{corollary}

\begin{theorem}[Consensus in \SOMminus\ cliques~\cite{aranda2024sound}]
\label{thm:som-consensus}
  Let $(G,\mathbf{B}^0,\mathbf{S}^0,\mu_G)$ be a \SOMminus\ model
  on an $n$-agent clique with $n\ge 3$.
  Then all agents converge to consensus.
\end{theorem}

% ===================================================================
\subsection{\texorpdfstring{\CBSOMminus}{CBSOM-}: Memoryless with Cognitive Biases}
\label{sec:cbsom-minus}

\begin{definition}[\CBSOMminus]\label{def:cbsom-minus}
  An SO model $(G,\mathbf{B}^0,\mathbf{S}^0,\mu_G)$ is
  \emph{memoryless with cognitive biases} (\CBSOMminus) if,
  for each $i\in A$ and $t\in\mathbb{N}$:
  \begin{align}
    B_i^{t+1} &= B_i^t
      + \sum_{j\in N(i)} I_{ji}\cdot S_j^t
        \cdot\beta_{i,j}(B_j^t - B_i^t),
    \label{eq:cbsom-minus-b}\\
    S_i^{t+1} &=
    \begin{cases}
      1 & \text{if }\;\bigl\lceil|N^t(i)|/2\bigr\rceil
          \le\bigl|\{j\in N^t(i)\mid
            B_i^t\sim_{\tau_i}B_j^t\}\bigr|,\\
      0 & \text{otherwise,}
    \end{cases}
    \label{eq:cbsom-minus-s}
  \end{align}
  where each $\beta_{i,j}:[-1,1]\to[-1,1]$ is a bias function.
\end{definition}

Equation~\eqref{eq:cbsom-minus-b} extends \eqref{eq:som-minus-b}
by replacing the linear disagreement term $(B_j^t-B_i^t)$ with the
biased term $\beta_{i,j}(B_j^t-B_i^t)$. The silence factor $S_j^t$
continues to zero out silent agents' contributions entirely; thus,
\emph{bias determines how vocal neighbors are processed, while
silence determines which neighbors are visible}. The two mechanisms
remain structurally decoupled. The following example illustrates both.

\begin{example}\label{ex:cbsom-minus}
  Consider a clique of three agents in which each agent gives weight
  $1/2$ to each of the other two; that is, $I_{ji}=1/2$ for $j\ne i$ and
  $I_{ii}=0$. Let $\mathbf{B}^0=(1,0.5,0)$, $\mathbf{S}^0=(1,1,1)$, and
  $\tau_i=0.4$ for every agent, and let every edge carry the bias
  $\beta_{i,j}(x)=x/2$, which halves each disagreement.
  By~\eqref{eq:cbsom-minus-b}, agent~1 updates its opinion to
  \[
    B_1^1 = 1+\tfrac{1}{2}\cdot\tfrac{1}{2}(0.5-1)
             +\tfrac{1}{2}\cdot\tfrac{1}{2}(0-1) = 0.625,
  \]
  and similarly $\mathbf{B}^1=(0.625,0.5,0.375)$. Without bias
  ($\beta_{i,j}(x)=x$) we would obtain $(0.25,0.5,0.75)$ instead: agents~1
  and~3 would swap sides, whereas the biased agents move only part of the
  way towards each other.

  Now consider silence. At $t=0$, every agent disagrees with both of its
  neighbors by more than $\tau_i=0.4$, so by~\eqref{eq:cbsom-minus-s} all
  agents fall silent: $\mathbf{S}^1=(0,0,0)$. With no vocal agent, no
  opinion changes, and $\mathbf{B}^2=\mathbf{B}^1$. At $t=1$ no agent hears
  any disagreement, so all agents speak again: $\mathbf{S}^2=(1,1,1)$. From
  then on all pairwise distances are at most $0.25<\tau_i$, so the agents
  remain vocal and approach the consensus value $0.5$. This episode of
  global silence is not a coincidence: Lemma~\ref{lem:no-perpetual-silence}
  in Section~\ref{sec:analysis} shows that global silence never lasts more
  than one round.
\end{example}

\begin{remark}[DeGroot as a special case]\label{rem:degroot-cbsom-minus}
  Setting $\tau_i=1$, $\mathbf{S}^0=\mathbf{1}_n$, and
  $\beta_{i,j}(x)=x$ recovers the classical DeGroot update from
  \CBSOMminus.
\end{remark}

% ===================================================================
\subsection{\texorpdfstring{\CBSOMplus}{CBSOM+}: Memory-Based with Cognitive Biases}
\label{sec:cbsom-plus}

\begin{definition}[\CBSOMplus]\label{def:cbsom-plus}
  An SO model $(G,\mathbf{B}^0,\mathbf{S}^0,\mu_G)$ is
  \emph{memory-based with cognitive biases} (\CBSOMplus) if
  $\mathbf{S}^0=\mathbf{1}_n$ and, for each $i\in A$ and
  $t\in\mathbb{N}$:
  \begin{align}
    B_i^{t+1} &= B_i^t
      + \sum_{j\in N(i)} I_{ji}
        \cdot\beta_{i,j}(\hat{B}_j^t - B_i^t),
    \label{eq:cbsom-plus-b}\\
    S_i^{t+1} &=
    \begin{cases}
      1 & \text{if }\;\bigl\lceil|N(i)|/2\bigr\rceil
          \le\bigl|\{j\in N(i)\mid
            B_i^t\sim_{\tau_i}\hat{B}_j^t\}\bigr|,\\
      0 & \text{otherwise,}
    \end{cases}
    \label{eq:cbsom-plus-s}
  \end{align}
  where $\hat{B}_j^t=B_j^{t_j'}$ with
  $t_j'=\max\{u\le t:S_j^u=1\}$.
\end{definition}

The model \CBSOMplus\ applies the bias function to historical public
opinions: $\beta_{i,j}$ filters how agent $i$ assimilates the last
expressed belief of neighbor~$j$, regardless of whether $j$ is currently
vocal. This creates a coupling between time and cognition that is absent
in \CBSOMminus. The frozen opinion of a silent agent is processed through
the bias of each of its neighbors at every step, until that agent speaks
again.

\begin{remark}[DeGroot as a special case]\label{rem:degroot-cbsom-plus}
  Setting $\tau_i=1$ and $\beta_{i,j}(x)=x$ recovers DeGroot from
  \CBSOMplus: perpetual vocality ensures $\hat{B}_j^t=B_j^t$ for
  all $j,t$, reducing \eqref{eq:cbsom-plus-b} to \eqref{eq:degroot}.
\end{remark}

% ===================================================================
\subsection{Relationship Among Models}
\label{sec:model-relations}

Table~\ref{tab:models} summarizes the model hierarchy.
The DeGroot model is a special case of both \CBSOMminus\ and \CBSOMplus\
(Remarks~\ref{rem:degroot-cbsom-minus}--\ref{rem:degroot-cbsom-plus}).
The \SOM\ models are the special cases of the \CBSOM\ models obtained
by setting $\beta_{i,j}(x)=x$. Therefore, every result proved for \CBSOM\
with identity biases recovers the corresponding \SOM\ result.

\begin{table}[!htbp]
\centering
\caption{Hierarchy of opinion models.
  $\checkmark$ = mechanism present; $-$ = absent or identity}
\label{tab:models}
\begin{tabular}{lccc}
  \hline
  \textbf{Model} & \textbf{Silence} & \textbf{Memory} & \textbf{Bias $\beta_{i,j}$}\\
  \hline
  DeGroot~\cite{DeGroot1974}               & $-$          & $-$  & identity\\
  \SOMminus~\cite{aranda2024sound}         & memoryless   & $-$  & identity\\
  \SOMplus~\cite{aranda2024sound}         & memory-based & $\checkmark$ & identity\\
  \CBSOMminus\ \emph{(this work)}          & memoryless   & $-$  & general\\
  \CBSOMplus\ \ \emph{(this work)}         & memory-based & $\checkmark$ & general\\
  \hline
\end{tabular}
\end{table}
